\documentclass[12pt]{article}
\usepackage[brazil]{babel}
\usepackage[utf8]{inputenc}
\usepackage[T1]{fontenc}
\usepackage{setspace}
\usepackage{geometry}
\usepackage{csquotes}

\geometry{a4paper, left=3cm, right=2cm, top=3cm, bottom=2cm}
\onehalfspacing

\title{EdTechs no Brasil e no cen\'ario internacional: revis\~ao de literatura (2014--2024)}
\author{Nome do(a) autor(a)}
\date{2026}

\begin{document}

\maketitle

\begin{abstract}
Nas \'ultimas d\'ecadas, empresas de tecnologia educacional (EdTechs) consolidaram-se como atores centrais em ecossistemas de inova\c{c}\~ao voltados \`a educa\c{c}\~ao b\'asica, ao ensino superior e \`a forma\c{c}\~ao ao longo da vida. O objetivo deste artigo \'e analisar, em perspectiva comparada, contribui\c{c}\~oes recentes da literatura sobre transforma\c{c}\~ao digital na educa\c{c}\~ao e sobre EdTechs no Brasil, em di\'alogo com tend\^encias e lacunas apontadas em revis\~oes internacionais publicadas entre 2014 e 2024. Para tanto, realizou-se uma revis\~ao narrativa de literatura, com inspira\c{c}\~ao em princ\'ipios de revis\~oes sistem\'aticas, privilegiando estudos de acesso aberto indexados em bases acad\^emicas e peri\'odicos cient\'ificos de educa\c{c}\~ao e tecnologia. Os resultados indicam converg\^encias quanto \`a centralidade de plataformas digitais, de dados e de intelig\^encia artificial (IA) nos processos de ensino, aprendizagem e gest\~ao, bem como assimetrias relevantes entre pa\'ises de alta renda e o contexto brasileiro. No Brasil, a literatura enfatiza desigualdades de conectividade, desafios de forma\c{c}\~ao docente e tens\~oes regulat\'orias na intera\c{c}\~ao entre redes p\'ublicas e startups, enquanto revis\~oes globais se concentram em experimentos com IA em ambientes bem equipados. Identificam-se lacunas relacionadas \`a avalia\c{c}\~ao de impacto de solu\c{c}\~oes EdTech e \`a representa\c{c}\~ao ainda limitada de pa\'ises de baixa e m\'edia renda em estudos de alcance internacional.

\textbf{Palavras-chave:} EdTechs; transforma\c{c}\~ao digital; intelig\^encia artificial em educa\c{c}\~ao; Brasil; revis\~ao de literatura.
\end{abstract}

\section{Introdu\c{c}\~ao}

A expans\~ao de tecnologias digitais de comunica\c{c}\~ao e informa\c{c}\~ao vem transformando, de modo acelerado, formas de ensinar, aprender e gerir sistemas educacionais em diferentes pa\'ises. Nesse cen\'ario, empresas de tecnologia educacional (EdTechs) deixam de ocupar lugar perif\'erico e passam a integrar de maneira estrutural os ecossistemas de inova\c{c}\~ao em educa\c{c}\~ao b\'asica, ensino superior e forma\c{c}\~ao profissional. No Brasil, esse movimento assume contornos particulares devido \`a combina\c{c}\~ao entre fortes desigualdades sociais, hist\'orico de infraestrutura digital prec\'aria e crescente demanda por solu\c{c}\~oes escal\'aveis que articulem qualidade pedag\'ogica e efici\^encia na gest\~ao.

A pandemia de COVID-19 funcionou como catalisador desse processo, impulsionando a ado\c{c}\~ao de plataformas digitais e de formatos de e-learning tanto em redes p\'ublicas de ensino b\'asico quanto em institui\c{c}\~oes de ensino superior (IES). Estudos recentes mostram que a chamada \emph{transforma\c{c}\~ao digital} (TD) em IES brasileiras n\~ao se limita \`a digitaliza\c{c}\~ao de conte\'udos, mas envolve reconfigura\c{c}\~oes em estrat\'egias institucionais, modelos de gest\~ao e rela\c{c}\~oes com novos atores, entre eles as EdTechs.

Em paralelo, revis\~oes sistem\'aticas internacionais documentam a expans\~ao de modelos de educa\c{c}\~ao aberta e a dist\^ancia (\emph{open and distance learning}, ODL) associados aos Objetivos de Desenvolvimento Sustent\'avel (ODS), bem como o avan\c{c}o da intelig\^encia artificial (IA) em contextos educacionais, com foco em personaliza\c{c}\~ao, anal\'iticas de aprendizagem e tutores inteligentes.[web:34][web:40] Este artigo busca integrar esses dois planos --- nacional e internacional --- de forma comparada, evidenciando converg\^encias, assimetrias, lacunas e oportunidades para o ecossistema brasileiro de EdTechs.

\section{Metodologia}

Adotou-se uma revis\~ao narrativa de literatura, com inspira\c{c}\~ao em princ\'ipios de revis\~oes sistem\'aticas, orientada por duas quest\~oes centrais: (i) como a literatura recente descreve a transforma\c{c}\~ao digital e o papel das EdTechs no sistema educacional brasileiro; e (ii) que tend\^encias e lacunas s\~ao apontadas por revis\~oes internacionais sobre tecnologias educacionais, em especial ODL e IA, entre 2014 e 2024.

A busca concentrou-se em estudos de acesso aberto em bases como SciELO, ERIC, Scopus e Web of Science, al\'em de peri\'odicos especializados em educa\c{c}\~ao e tecnologia. Foram considerados artigos publicados entre 2014 e 2024, com foco em pelo menos um dos seguintes eixos: (a) transforma\c{c}\~ao digital em IES brasileiras; (b) estrat\'egias digitais de redes de ensino durante a pandemia de COVID-19; (c) revis\~oes sistem\'aticas internacionais sobre educa\c{c}\~ao aberta e a dist\^ancia associada aos ODS; e (d) revis\~oes sistem\'aticas sobre IA na educa\c{c}\~ao.

Para o eixo brasileiro, privilegiaram-se estudos que abordam TD em IES a partir de referenciais estrat\'egicos, como a Vis\~ao Baseada em Recursos (\emph{Resource-Based View} -- RBV), e investiga\c{c}\~oes sobre estrat\'egias digitais durante a pandemia, com aten\c{c}\~ao para desigualdades socioeducacionais e conectividade.[page:1][web:22] No eixo internacional, foram selecionadas revis\~oes sistem\'aticas que sintetizam experi\^encias de ODL relacionadas ao ODS 4 (educa\c{c}\~ao de qualidade) e estudos sobre IA na educa\c{c}\~ao, com \^enfase em tend\^encias globais e lacunas metodol\'ogicas.[web:34][web:40]

\section{EdTechs e transforma\c{c}\~ao digital no Brasil}

A literatura brasileira recente enfatiza que a TD em IES foi intensificada pela pandemia, mas n\~ao se reduz a uma resposta conjuntural. Jacociunas, Verschoore e Monticelli prop\~oem um \emph{framework} estrat\'egico para TD em IES brasileiras, fundamentado na RBV, que identifica oito atributos cr\'iticos: efetividade dos cursos, qualidade dos programas educacionais, capacita\c{c}\~ao de professores para o e-learning, sistema de gest\~ao apropriado, processo de vendas ativo, engajamento dos profissionais, tecnologia e desenvolvimento \'agil.[page:1] Esses atributos, articulados a uma cultura empreendedora, condicionam a capacidade de as institui\c{c}\~oes competirem em um cen\'ario no qual EdTechs especializados disputam estudantes e oferecimento de servi\c{c}os para o setor educacional.

Os autores destacam que a simples aquisi\c{c}\~ao de tecnologias n\~ao garante vantagem competitiva, sendo necess\'aria uma reconfigura\c{c}\~ao mais ampla de processos e de cultura organizacional, incluindo estruturas de decis\~ao menos burocr\'aticas e maior abertura \`a inova\c{c}\~ao.[page:1] Nesse sentido, parcerias com EdTechs --- provendo plataformas de gest\~ao, ambientes virtuais de aprendizagem, solu\c{c}\~oes de anal\'iticas e servi\c{c}os de suporte docente --- aparecem como estrat\'egias relevantes para que IES ampliem capacidade tecnol\'ogica sem necessariamente internalizar todas as compet\^encias.

No campo da educa\c{c}\~ao b\'asica, estudos sobre estrat\'egias digitais durante a pandemia mostram que, diante da aus\^encia de diretrizes nacionais unificadas, estados e munic\'ipios desenharam solu\c{c}\~oes pr\'oprias de ensino remoto, com grande varia\c{c}\~ao de qualidade e de acesso.[web:22] Evid\^encias apontam que desigualdades pr\'evias de conectividade, disponibilidade de dispositivos e apoio familiar se converteram em lacunas de aprendizagem ainda mais profundas, especialmente entre estudantes de redes p\'ublicas em regi\~oes rurais ou periferias urbanas.[web:22]

Nesse contexto, EdTechs que oferecem solu\c{c}\~oes voltadas \`a conectividade escolar, gest\~ao em nuvem, recursos multim\'idia de baixo consumo de dados e apoio ao trabalho docente em contextos vulner\'aveis surgem como atores estrat\'egicos. No entanto, a literatura destaca barreiras importantes para a contrata\c{c}\~ao p\'ublica e a escala dessas iniciativas, incluindo processos licitat\'orios complexos, inseguran\c{c}a jur\'idica e dificuldades de integra\c{c}\~ao de dados entre sistemas diversos.[web:18][web:22]

\section{Cen\'ario internacional}

No plano internacional, revis\~oes sistem\'aticas sobre educa\c{c}\~ao aberta e a dist\^ancia evidenciam que, entre 2014 e 2024, houve aumento significativo do n\'umero de estudos que associam ODL ao alcance do ODS 4, relativo \`a educa\c{c}\~ao de qualidade.[web:34] Essas revis\~oes destacam o papel de plataformas digitais, cursos online massivos e ambientes virtuais de aprendizagem na amplia\c{c}\~ao do acesso, na promo\c{c}\~ao da aprendizagem ao longo da vida e na internacionaliza\c{c}\~ao de experi\^encias formativas, sobretudo em pa\'ises de alta renda com infraestrutura consolidada de educa\c{c}\~ao online.[web:34]

Em paralelo, revis\~oes sobre IA na educa\c{c}\~ao, com destaque para estudos voltados \`a matem\'atica, mostram crescimento acentuado de publica\c{c}\~oes a partir de 2020, com concentra\c{c}\~ao em pa\'ises de alta renda e em contextos de ensino m\'edio e superior.[web:40] As pesquisas abordam desde \emph{chatbots} baseados em grandes modelos de linguagem at\'e sistemas adaptativos e plataformas de anal\'iticas de aprendizagem, sendo organizadas em eixos como apoio ao pensamento matem\'atico, personaliza\c{c}\~ao da aprendizagem, desenvolvimento profissional docente, quest\~oes \'eticas e de equidade e inova\c{c}\~oes multimodais com rob\'otica.[web:40][web:41]

Apesar desse avan\c{c}o, as revis\~oes internacionais ressaltam limita\c{c}\~oes importantes. Entre elas est\~ao a predomin\^ancia de estudos qualitativos ou descritivos de pequena escala, a escassez de investiga\c{c}\~oes longitudinais que acompanhem impactos de longo prazo e a baixa representa\c{c}\~ao de pa\'ises de baixa e m\'edia renda em amostras globais.[web:34][web:40] Tais lacunas indicam que, embora a fronteira tecnol\'ogica avance rapidamente em contextos de alta renda, ainda existe pouca evid\^encia robusta sobre efetividade, equidade e implica\c{c}\~oes \'eticas do uso intensivo de IA e de plataformas EdTech em sistemas educacionais marcados por desigualdades estruturais.

\section{Brasil e cen\'ario internacional: converg\^encias e assimetrias}

Comparando os achados da literatura brasileira com as revis\~oes internacionais, observam-se converg\^encias e assimetrias. Entre as converg\^encias, destaca-se o reconhecimento de que plataformas digitais, dados e IA passam a integrar de forma indissoci\'avel o cotidiano de processos de ensino, aprendizagem e gest\~ao, tanto em IES quanto em redes de educa\c{c}\~ao b\'asica.[page:1][web:34][web:40] Tanto no Brasil quanto em pa\'ises de alta renda, TD \'e vista como processo multidimensional, que envolve dimens\~oes tecnol\'ogicas, organizacionais, pedag\'ogicas e culturais.

No entanto, as assimetrias s\~ao expressivas. Enquanto revis\~oes globais tendem a enfatizar experimentos com IA e ODL em ambientes bem equipados, com foco em efici\^encia, personaliza\c{c}\~ao e inova\c{c}\~ao pedag\'ogica, a literatura brasileira destaca de forma mais contundente desafios de conectividade, forma\c{c}\~ao docente e regula\c{c}\~ao na intera\c{c}\~ao entre setor p\'ublico e EdTechs.[web:18][web:22][web:34] Ademais, h\'a forte preocupa\c{c}\~ao com a sustentabilidade econ\^omica das startups e com o risco de aprofundamento de desigualdades caso solu\c{c}\~oes tecnol\'ogicas sejam desenhadas sem considera\c{c}\~ao das condi\c{c}\~oes reais de escolas e comunidades.

Outra assimetria importante diz respeito \`a presen\c{c}a de pa\'ises como o Brasil nas pr\'oprias revis\~oes internacionais. Mesmo quando estudos mencionam Am\'erica Latina, a quantidade de trabalhos emp\'iricos em contextos de baixa renda ainda \'e reduzida, o que dificulta generaliza\c{c}\~oes sobre impacto e equidade das solu\c{c}\~oes EdTech nesses ambientes.[web:34][web:40] Isso refor\c{c}a a import\^ancia de fortalecer a produ\c{c}\~ao acad\^emica local, capaz de analisar e documentar experi\^encias com tecnologias educacionais em diferentes regi\~oes e segmentos do sistema educacional brasileiro.

\section{Lacunas e oportunidades}

A partir da literatura analisada, podem ser identificadas lacunas e oportunidades de pesquisa e de atua\c{c}\~ao para o ecossistema de EdTechs no Brasil. Em primeiro lugar, estudos nacionais indicam a necessidade de avalia\c{c}\~oes de impacto mais robustas, que acompanhem, em m\'edio e longo prazo, efeitos de solu\c{c}\~oes tecnologicamente intensivas sobre aprendizagem, perman\^encia e equidade, tanto em educa\c{c}\~ao b\'asica quanto em educa\c{c}\~ao superior.[page:1][web:22] Tais avalia\c{c}\~oes podem combinar m\'etodos quantitativos e qualitativos, articulando indicadores de desempenho com narrativas de professores, estudantes e gestores.

Em segundo lugar, h\'a oportunidade de aprofundar estudos sobre implementa\c{c}\~ao de \emph{frameworks} de TD, como o proposto por Jacociunas, Verschoore e Monticelli, em diferentes tipos de IES (p\'ublicas, privadas, presenciais e a dist\^ancia), mapeando condi\c{c}\~oes para sua institucionaliza\c{c}\~ao efetiva e os resultados decorrentes.[page:1] Esse tipo de pesquisa pode gerar evid\^encias valiosas para decisores institucionais e formuladores de pol\'itica.

Em terceiro lugar, as revis\~oes internacionais sobre ODL e IA indicam necessidade de ampliar a diversidade geogr\'afica e socioecon\^omica de amostras, incluindo mais estudos em pa\'ises de baixa e m\'edia renda e em regi\~oes com infraestrutura limitada.[web:34][web:40] Para o Brasil, isso se traduz na oportunidade de desenvolver e investigar solu\c{c}\~oes EdTech orientadas a contextos de baixa conectividade, escolas rurais, popula\c{c}\~oes ind\'igenas e comunidades perif\'ericas, articulando tecnologias emergentes a abordagens pedag\'ogicas inclusivas e culturalmente sens\'iveis.

Por fim, tanto a produ\c{c}\~ao nacional quanto as revis\~oes globais convergem na mensagem de que a maturidade do setor de EdTechs n\~ao dever\'a ser medida apenas pelo n\'umero de startups ou pelo volume de investimento, mas pela capacidade de produzir evid\^encias de impacto educacional, pela integra\c{c}\~ao com pol\'iticas p\'ublicas e pela aten\c{c}\~ao \`as dimens\~oes \'eticas, de prote\c{c}\~ao de dados e de justi\c{c}a social no desenho e na ado\c{c}\~ao de tecnologias.[page:1][web:34][web:40]

\section{Considera\c{c}\~oes finais}

Este artigo apresentou uma revis\~ao narrativa de literatura sobre EdTechs e transforma\c{c}\~ao digital na educa\c{c}\~ao, com foco no contexto brasileiro e em di\'alogo com revis\~oes internacionais publicadas entre 2014 e 2024. A an\'alise mostrou que, embora haja converg\^encias quanto \`a centralidade de plataformas digitais, dados e IA em processos educacionais, persistem assimetrias importantes entre pa\'ises de alta renda e contextos como o Brasil, marcados por desigualdades estruturais.

No Brasil, estudos destacam a import\^ancia de articular recursos institucionais hist\'oricos a capacidades digitais emergentes, bem como de enfrentar desafios de conectividade, forma\c{c}\~ao docente e regula\c{c}\~ao na intera\c{c}\~ao com EdTechs.[page:1][web:22] No plano internacional, revis\~oes sobre ODL e IA evidenciam tanto o potencial de tecnologias educacionais para ampliar acesso e personalizar experi\^encias quanto lacunas referentes \`a avalia\c{c}\~ao de impacto e \`a representatividade de pa\'ises de baixa e m\'edia renda.[web:34][web:40]

Como agenda de pesquisa e de desenvolvimento, destaca-se a necessidade de produzir evid\^encias de impacto de solu\c{c}\~oes EdTech em contextos reais de aprendizagem, de fortalecer parcerias entre IES, redes p\'ublicas e startups e de orientar a inova\c{c}\~ao tecnol\'ogica por princ\'ipios de equidade, inclus\~ao e justi\c{c}a social. Tais dire\c{c}\~oes s\~ao fundamentais para que o ecossistema brasileiro de EdTechs contribua de forma efetiva para o direito \`a educa\c{c}\~ao de qualidade para todas e todos.

\section*{Refer\^encias}

JACOCIUNAS, T.; VERSCHOORE, J.; MONTICELLI, J. Transforma\c{c}\~ao digital de institui\c{c}\~oes de ensino superior: um framework para a tomada de decis\~ao estrat\'egica. \textit{Revista Internacional de Educa\c{c}\~ao Superior}, v. 10, e024036, 2024.[page:1]

BRAZILIAN digital learning strategies during the COVID-19 pandemic and their impact on socio-educational gaps. \textit{Current Issues in Comparative Education}, v. 24, n. 1, p. 87--102, 2022.[web:22]

TRANSFORMATION digital de institui\c{c}\~oes de ensino superior. \textit{Revista Internacional de Educa\c{c}\~ao Superior}, v. 10, e024036, 2024.[web:18]

SYSTEMATIC Review of the Role of Open and Distance Education in Achieving the Sustainable Development Goals (SDGs). \textit{International Journal of Learning, Teaching and Educational Research}, v. 23, n. 12, 2024.[web:34]

ARTIFICIAL Intelligence in Mathematics Education: A Systematic Review of Global Trends and Emerging Themes. \textit{Journal of Mathematics, Science and Technology Education}, 2025.[web:40]

TRENDS and Security Challenges of Artificial Intelligence Use in Higher Education: A Systematic Review. \textit{IEEE Transactions on Education}, 2025.[web:41]

\end{document}
